% -------------------------------------------------------
% Informationen und Einstellungen für Ihre Abschlussarbeit
%

\newcommand{\seesettingsmarker}

% Sprache für das Dokument festlegen
\newcommand{\hsmasprache}{de} %de für Deutsch oder en für Englisch


% Abgabeform festlegen
% Bei einer digitalen Abgabe, wird das Dokument einseitig erzeugt und der Titel wird
% zentriert.
\newcommand{\hsmaabgabe}{digital} % Abgabe erfolgt für Fakultät I digital. Optionen hier sind für anderen Fakultäten: "papier" oder "digital".


% Flags für Veröffentlichung, Sperrvermerk
\newcommand{\hsmapublizieren}{stud}
% Wird einer Veröffentlichung zugestimmt?
% Optionen:
% opensource = Druck der CC Lizenz mit By SA (Standard)
% hs = Veröffentlichung an der Technischen Hochschule und auf Hochschulservern
% stud = kein opensource und keine veröffentlichung auf den Hochschulservern
% vertraulich = Arbeit darf nicht veröffentlicht werden und erhält einen Sperrvermerk (Nur nach Absprache mit Betreuer setzen!)


\newcommand{\genderhinweis}{nogender}     % Soll der Gender-Hinweis angezeigt werden? ja=gender, nein = nogender; Genderhinweis wird nur in deutscher Sprache angezeigt!


\newcommand{\hsmaquellcode}{sourcecode} % Verwenden Sie Quellcode in Ihrer Arbeit? ja=sourcecode, nein= nosourcecode

\newcommand{\hsmasymbole}{nosymbole} % Verwenden Sie viele Symbole in Ihrer Arbeit, welche in einem Symbolverzeichnis aufgeführt werden sollen? ja=symbole, nein= nosymbole


\newcommand{\hsmaglossar}{noglossar} % Verwenden Sie Begriffserklärungen nicht Abkürzungen in Ihrer Arbeit? ja=glossar, nein= noglossar

\newcommand{\hsmatc}{notc} % Verwenden der Änderungsmarkierung. Änderungsmarkierung aktiv und eine Liste der Änderungen wird angezeigt = tc, Keine Änderungsmarkierung und keine Ausgabe der Änderungen = notc




% Titel der Arbeit auf Deutsch
\newcommand{\hsmatitelde}{Echtzeitkommunikation mit WebSockets und React: Implementierung einer Spiellobby mit Live-Benachrichtigungen und gemeinsamem Lobby-Zustand}

% Titel der Arbeit auf Englisch
\newcommand{\hsmatitelen}{Real-time communication with WebSockets and React: implementation of a game lobby with live notifications and shared lobby state}

% Weitere Informationen zur Arbeit
\newcommand{\hsmaort}{Mannheim}          % Ort
\newcommand{\hsmaautorvname}{David}        % Vorname(n)
\newcommand{\hsmaautornname}{Groys} % Nachname(n)
\newcommand{\hsmaabgabedatum}{2026-12-15}% TODO: echtes Abgabedatum% Datum der Abgabe in dem Format JJJJ-MM-TT

\newcommand{\hsmafirma}{} % Firma bei der die Arbeit durchgeführt wurde
\newcommand{\hsmabetreuer}{} % Betreuer an der Hochschule
\newcommand{\hsmazweitkorrektor}{}   % Betreuer im Unternehmen oder Zweitkorrektor

\newcommand{\hsmafakultaet}{I}    % I für Informatik oder E, S, B, D, M, N, W, V
\newcommand{\hsmastudiengang}{IB} % IB IMB UIB CSB IM MTB (weitere siehe titleblatt.tex)


% -------------------------------------------------------
% Ergänzungen für eine Seminararbeit (statt Bachelor-Thesis)
\newcommand{\hsmaarbeitde}{Seminararbeit}
\newcommand{\hsmaarbeiten}{Seminar Paper}
\newcommand{\hsmamodul}{Ang + NEN}

% Autoren (Titelseite und Kennzeichnung der Abschnitte)
\newcommand{\autorDG}{David Groys}
\newcommand{\autorSE}{Selim Esser}
\newcommand{\hsmaautorentitel}{\autorDG\ (3012642)\\[2mm]\autorSE\ (2211482)}
\newcommand{\hsmaautorenkurz}{David Groys, Selim Esser}
